\documentclass[../main.tex]{subfiles}
\begin{document}

\section{Phương pháp xây dựng luật phát hiện S7comm}
\label{sec:detection-overview}

Chương này trình bày thiết kế tập luật Suricata phát hiện hành vi tấn công trên giao thức Siemens S7comm, căn cứ phân tích lưu lượng ở Chương~\ref{sec:attack-scenarios}. Phương pháp thuộc nhóm \emph{phát hiện dựa trên chữ ký}: mỗi luật khớp tổ hợp byte cố định trong payload TCP/102 (TPKT, COTP, S7 Header, mã chức năng). 

Toàn bộ luật nằm tại \texttt{detect/rules/s7comm.rules}. Mã \texttt{sid} phân theo nhóm: \texttt{100000x} (thiết lập phiên, trinh sát), \texttt{100001x} (đọc ghi biến), \texttt{100002x} (download upload block), \texttt{100003x} (điều khiển Start/Stop).

\section{Cơ sở chữ ký trên payload S7comm}
\label{sec:rule-foundation}

S7comm trên TCP/102 có cấu trúc TCP $\rightarrow$ TPKT $\rightarrow$ COTP $\rightarrow$ S7comm. Bảng~\ref{tab:s7-payload-offsets} liệt kê offset byte dùng trong luật.

\begin{table}[H]
    \centering
    \caption{Offset byte trong payload TCP/102 (S7comm Data TPDU)}
    \label{tab:s7-payload-offsets}
    \begin{tabular}{|c|c|p{7.2cm}|}
        \hline
        \textbf{Offset} & \textbf{Giá trị} & \textbf{Trường} \\
        \hline
        0--1 & \texttt{03 00} & TPKT \\
        \hline
        4--6 & \texttt{02 F0 80} & COTP Data TPDU \\
        \hline
        7 & \texttt{32} & S7 protocol id \\
        \hline
        8 & \texttt{01}/\texttt{03}/\texttt{07} & \texttt{ROSCTR}: Job / Ack\_Data / Userdata \\
        \hline
        17 & --- & Function code (Job, Userdata) \\
        \hline
        19 & --- & Function code (Ack\_Data) \\
        \hline
    \end{tabular}
\end{table}

Luật chiều client $\rightarrow$ PLC dùng \texttt{flow:to\_server,established} và \texttt{|02 F0 80 32 01|} tại offset~4 (\texttt{ROSCTR=Job}). Luật Userdata dùng \texttt{|02 F0 80 32 07|}. Luật chiều PLC $\rightarrow$ client dùng \texttt{flow:from\_server,established} và \texttt{|02 F0 80 32 03|} (\texttt{ROSCTR=Ack\_Data}); function code tại offset~19.

\section{Luật theo kịch bản tấn công}
\label{sec:rules-by-scenario}

\subsection{Reconnaissance}
\label{sec:rules-recon}

Căn cứ Mục~\ref{sec:reconnaissance}: quét TCP/102 (Nmap \texttt{s7-info}), đọc SZL (\texttt{0x0011}, \texttt{0x001c}), liệt kê block (Snap7).

\paragraph{Luật \texttt{sid:1000001}.}
\begin{footnotesize}
\begin{verbatim}
alert tcp any any -> $HOME_NET 102 (msg:"ICS S7COMM Setup Communication to PLC";
  flow:to_server,established; content:"|03 00|"; offset:0; depth:2;
  content:"|02 F0 80 32 01|"; offset:4; depth:5; content:"|F0|";
  offset:17; depth:1; sid:1000001; rev:1;
  metadata:service iso-tsap, protocol s7comm, attack reconnaissance; priority:3;)
\end{verbatim}
\end{footnotesize}

\noindent\textbf{Phân tích.}
Khớp Job (\texttt{ROSCTR=0x01}) với function \texttt{0xf0} (Setup Communication)---bước bắt buộc trước mọi thao tác S7comm trong các kịch bản Chương~3. \texttt{sid:1000001} không đủ để kết luận tấn công; dùng làm dấu hiệu mở phiên hoặc tương quan với các \texttt{sid} trinh sát tiếp theo.

\paragraph{Luật \texttt{sid:1000002}.}
\begin{footnotesize}
\begin{verbatim}
alert tcp any any -> $HOME_NET 102 (msg:"ICS S7COMM Read SZL request - PLC
  information enumeration"; flow:to_server,established; content:"|03 00|";
  offset:0; depth:2; content:"|02 F0 80 32 07|"; offset:4; depth:5;
  content:"|00 01 12 04 11 44 01 00|"; offset:17; depth:8; sid:1000002; rev:1;
  metadata:service iso-tsap, protocol s7comm, attack reconnaissance; priority:2;)
\end{verbatim}
\end{footnotesize}

\noindent\textbf{Phân tích.}
\texttt{ROSCTR=0x07} (Userdata) và mẫu \texttt{|00 01 12 04 11 44 01 00|} xác định nhóm CPU functions, subfunction Read SZL. \texttt{sid:1000002} phát hiện mọi yêu cầu đọc SZL; các \texttt{sid:1000003} và \texttt{sid:1000004} tinh chỉnh theo SZL-ID quan sát trong PCAP reconnaissance.

\paragraph{Luật \texttt{sid:1000003}.}
\begin{footnotesize}
\begin{verbatim}
alert tcp any any -> $HOME_NET 102 (msg:"ICS S7COMM Read SZL module identification
  0x0011"; flow:to_server,established; content:"|03 00|"; offset:0; depth:2;
  content:"|02 F0 80 32 07|"; offset:4; depth:5;
  content:"|00 01 12 04 11 44 01 00|"; offset:17; depth:8;
  content:"|00 11 00 01|"; distance:0; within:12; sid:1000003; rev:1;
  metadata:service iso-tsap, protocol s7comm, attack reconnaissance; priority:2;)
\end{verbatim}
\end{footnotesize}

\noindent\textbf{Phân tích.}
Bổ sung \texttt{|00 11 00 01|} (SZL-ID \texttt{0x0011}, index \texttt{0x0001})---trường module identification (order number CPU) trong output Nmap \texttt{s7-info}. \texttt{sid:1000003} chứng minh fingerprint PLC cụ thể hơn \texttt{sid:1000002}.

\paragraph{Luật \texttt{sid:1000004}.}
\begin{footnotesize}
\begin{verbatim}
alert tcp any any -> $HOME_NET 102 (msg:"ICS S7COMM Read SZL communication status
  0x001c"; flow:to_server,established; content:"|03 00|"; offset:0; depth:2;
  content:"|02 F0 80 32 07|"; offset:4; depth:5;
  content:"|00 01 12 04 11 44 01 00|"; offset:17; depth:8;
  content:"|00 1C 00 00|"; distance:0; within:12; sid:1000004; rev:1;
  metadata:service iso-tsap, protocol s7comm, attack reconnaissance; priority:2;)
\end{verbatim}
\end{footnotesize}

\noindent\textbf{Phân tích.}
Khớp SZL-ID \texttt{0x001c} (communication status): serial number, system name, module type---khớp bước thu thập metadata PLC trong reconnaissance.

\paragraph{Luật \texttt{sid:1000005}.}
\begin{footnotesize}
\begin{verbatim}
alert tcp any any -> $HOME_NET 102 (msg:"ICS S7COMM List Blocks request - PLC block
  inventory enumeration"; flow:to_server,established; content:"|03 00|";
  offset:0; depth:2; content:"|02 F0 80 32 07|"; offset:4; depth:5;
  content:"|00 01 12 04 11 13 01 00|"; offset:17; depth:8; sid:1000005; rev:1;
  metadata:service iso-tsap, protocol s7comm, attack reconnaissance; priority:2;)
\end{verbatim}
\end{footnotesize}

\noindent\textbf{Phân tích.}
Mẫu \texttt{|00 01 12 04 11 13 01 00|}: function group Block functions, subfunction List blocks. Tương ứng \texttt{scanBlock.py} (\texttt{.list\_blocks()}). Đánh dấu giai đoạn kiểm kê OB/DB/FC/FB trước upload/download hoặc khai thác sâu.

\subsection{Phát lại lệnh Start/Stop PLC}
\label{sec:rules-start-stop}

Căn cứ Mục~\ref{sec:start_stop_plc}: replay payload \texttt{PLC Control 0x28} và \texttt{PLC Stop 0x29}, PI-Service \texttt{P\_PROGRAM}.

\paragraph{Luật \texttt{sid:1000030}.}
\begin{footnotesize}
\begin{verbatim}
alert tcp any any -> $HOME_NET 102 (msg:"ICS S7COMM PLC Control command";
  flow:to_server,established; content:"|03 00|"; offset:0; depth:2;
  content:"|02 F0 80 32 01|"; offset:4; depth:5; content:"|28|";
  offset:17; depth:1; sid:1000030; rev:1;
  metadata:service iso-tsap, protocol s7comm, attack plc_control; priority:1;)
\end{verbatim}
\end{footnotesize}

\noindent\textbf{Phân tích.}
Phát hiện mọi lệnh PLC Control (\texttt{0x28}): Start CPU, activate block, copy RAM to ROM, \ldots. \texttt{priority:1} phản ánh mức độ nghiêm trọng---ảnh hưởng trực tiếp trạng thái vận hành PLC.

\paragraph{Luật \texttt{sid:1000031}.}
\begin{footnotesize}
\begin{verbatim}
alert tcp any any -> $HOME_NET 102 (msg:"ICS S7COMM replayed Start CPU P_PROGRAM";
  flow:to_server,established; content:"|03 00|"; offset:0; depth:2;
  content:"|02 F0 80 32 01|"; offset:4; depth:5; content:"|28|";
  offset:17; depth:1; content:"P_PROGRAM"; distance:0; within:24;
  sid:1000031; rev:1;
  metadata:service iso-tsap, protocol s7comm, attack start_stop_replay; priority:1;)
\end{verbatim}
\end{footnotesize}

\noindent\textbf{Phân tích.}
Bổ sung chuỗi ASCII \texttt{P\_PROGRAM} trong 24 byte sau function code---khớp payload replay Start CPU từ PCAP TIA Portal. \texttt{sid:1000031} chuyên biệt hơn \texttt{sid:1000030} cho kịch bản replay đồ án.

\paragraph{Luật \texttt{sid:1000032}.}
\begin{footnotesize}
\begin{verbatim}
alert tcp any any -> $HOME_NET 102 (msg:"ICS S7COMM PLC Stop command";
  flow:to_server,established; content:"|03 00|"; offset:0; depth:2;
  content:"|02 F0 80 32 01|"; offset:4; depth:5; content:"|29|";
  offset:17; depth:1; sid:1000032; rev:1;
  metadata:service iso-tsap, protocol s7comm, attack plc_stop; priority:1;)
\end{verbatim}
\end{footnotesize}

\noindent\textbf{Phân tích.}
Khớp function \texttt{0x29} (PLC Stop). Tương đương \texttt{sid:1000030} nhưng cho lệnh dừng CPU.

\paragraph{Luật \texttt{sid:1000033}.}
\begin{footnotesize}
\begin{verbatim}
alert tcp any any -> $HOME_NET 102 (msg:"ICS S7COMM replayed Stop CPU P_PROGRAM";
  flow:to_server,established; content:"|03 00|"; offset:0; depth:2;
  content:"|02 F0 80 32 01|"; offset:4; depth:5; content:"|29|";
  offset:17; depth:1; content:"P_PROGRAM"; distance:0; within:24;
  sid:1000033; rev:1;
  metadata:service iso-tsap, protocol s7comm, attack start_stop_replay; priority:1;)
\end{verbatim}
\end{footnotesize}

\noindent\textbf{Phân tích.}
Đối xứng \texttt{sid:1000031}: Stop CPU với PI-Service \texttt{P\_PROGRAM}. Khớp payload replay Stop trong lab.

\subsection{Upload và Download block}
\label{sec:rules-up-download}

Căn cứ Mục~\ref{sec:down_up_program}: chuỗi function \texttt{0x1a}--\texttt{0x1f} trên traffic Siemens hoặc mô phỏng Wireshark.

\paragraph{Luật \texttt{sid:1000020}.}
\begin{footnotesize}
\begin{verbatim}
alert tcp any any -> $HOME_NET 102 (msg:"ICS S7COMM Request Download block to PLC";
  flow:to_server,established; content:"|03 00|"; offset:0; depth:2;
  content:"|02 F0 80 32 01|"; offset:4; depth:5; content:"|1A|";
  offset:17; depth:1; sid:1000020; rev:1;
  metadata:service iso-tsap, protocol s7comm, attack program_download; priority:1;)
\end{verbatim}
\end{footnotesize}

\noindent\textbf{Phân tích.}
Khớp Request Download (\texttt{0x1a}): client yêu cầu PLC chuẩn bị nhận block chương trình.

\paragraph{Luật \texttt{sid:1000021}.}
\begin{footnotesize}
\begin{verbatim}
alert tcp $HOME_NET 102 -> any any (msg:"ICS S7COMM PLC requests Download Block data";
  flow:from_server,established; content:"|03 00|"; offset:0; depth:2;
  content:"|02 F0 80 32 01|"; offset:4; depth:5; content:"|1B|";
  offset:17; depth:1; sid:1000021; rev:1;
  metadata:service iso-tsap, protocol s7comm, attack program_download; priority:1;)
\end{verbatim}
\end{footnotesize}

\noindent\textbf{Phân tích.}
Chiều PLC $\rightarrow$ client; Job \texttt{0x1b}: PLC chủ động yêu cầu block dữ liệu download. Khác \texttt{sid:1000026} về hướng và \texttt{ROSCTR}.

\paragraph{Luật \texttt{sid:1000026}.}
\begin{footnotesize}
\begin{verbatim}
alert tcp any any -> $HOME_NET 102 (msg:"ICS S7COMM client sends Download Block data
  to PLC"; flow:to_server,established; content:"|03 00|"; offset:0; depth:2;
  content:"|02 F0 80 32 03|"; offset:4; depth:5; content:"|1B|";
  offset:19; depth:1; sid:1000026; rev:1;
  metadata:service iso-tsap, protocol s7comm, attack program_download; priority:1;)
\end{verbatim}
\end{footnotesize}

\noindent\textbf{Phân tích.}
Ack\_Data (\texttt{ROSCTR=0x03}), function \texttt{0x1b} tại offset~19: client gửi nội dung block xuống PLC. Cặp \texttt{sid:1000021}/\texttt{sid:1000026} mô tả hai chiều của giai đoạn Download Block.

\paragraph{Luật \texttt{sid:1000022}.}
\begin{footnotesize}
\begin{verbatim}
alert tcp any any -> $HOME_NET 102 (msg:"ICS S7COMM Download Ended";
  flow:to_server,established; content:"|03 00|"; offset:0; depth:2;
  content:"|02 F0 80 32 01|"; offset:4; depth:5; content:"|1C|";
  offset:17; depth:1; sid:1000022; rev:1;
  metadata:service iso-tsap, protocol s7comm, attack program_download; priority:1;)
\end{verbatim}
\end{footnotesize}

\noindent\textbf{Phân tích.}
Kết thúc chuỗi download (\texttt{0x1c}).

\paragraph{Luật \texttt{sid:1000023}.}
\begin{footnotesize}
\begin{verbatim}
alert tcp any any -> $HOME_NET 102 (msg:"ICS S7COMM Start Upload block from PLC";
  flow:to_server,established; content:"|03 00|"; offset:0; depth:2;
  content:"|02 F0 80 32 01|"; offset:4; depth:5; content:"|1D|";
  offset:17; depth:1; sid:1000023; rev:1;
  metadata:service iso-tsap, protocol s7comm, attack program_upload; priority:1;)
\end{verbatim}
\end{footnotesize}

\noindent\textbf{Phân tích.}
Start Upload (\texttt{0x1d}): client yêu cầu lấy chương trình từ PLC.

\paragraph{Luật \texttt{sid:1000024}.}
\begin{footnotesize}
\begin{verbatim}
alert tcp any any -> $HOME_NET 102 (msg:"ICS S7COMM Upload Block request";
  flow:to_server,established; content:"|03 00|"; offset:0; depth:2;
  content:"|02 F0 80 32 01|"; offset:4; depth:5; content:"|1E|";
  offset:17; depth:1; sid:1000024; rev:1;
  metadata:service iso-tsap, protocol s7comm, attack program_upload; priority:1;)
\end{verbatim}
\end{footnotesize}

\noindent\textbf{Phân tích.}
Upload Block (\texttt{0x1e}): truyền từng block dữ liệu từ PLC về client.

\paragraph{Luật \texttt{sid:1000025}.}
\begin{footnotesize}
\begin{verbatim}
alert tcp any any -> $HOME_NET 102 (msg:"ICS S7COMM End Upload";
  flow:to_server,established; content:"|03 00|"; offset:0; depth:2;
  content:"|02 F0 80 32 01|"; offset:4; depth:5; content:"|1F|";
  offset:17; depth:1; sid:1000025; rev:1;
  metadata:service iso-tsap, protocol s7comm, attack program_upload; priority:1;)
\end{verbatim}
\end{footnotesize}

\noindent\textbf{Phân tích.}
End Upload (\texttt{0x1f}): kết thúc chuỗi upload.

\subsection{Ghi biến quá trình và quan sát MITM}
\label{sec:rules-stuxnet}

Căn cứ Mục~\ref{sec:stuxnet_mitm_sim}: ghi \texttt{DB1.DBW2} (\texttt{iSetpoint}) \texttt{0x04b0}$\rightarrow$\texttt{0x07d0}; can thiệp \texttt{Ack\_Data} phía HMI.

\paragraph{Luật \texttt{sid:1000010}.}
\begin{footnotesize}
\begin{verbatim}
alert tcp any any -> $HOME_NET 102 (msg:"ICS S7COMM Write Var request to PLC";
  flow:to_server,established; content:"|03 00|"; offset:0; depth:2;
  content:"|02 F0 80 32 01|"; offset:4; depth:5; content:"|05|";
  offset:17; depth:1; sid:1000010; rev:1;
  metadata:service iso-tsap, protocol s7comm, attack modify_process_variable;
  priority:1;)
\end{verbatim}
\end{footnotesize}

\noindent\textbf{Phân tích.}
Function \texttt{0x05} (Write Var). Cảnh báo mọi ghi biến vào PLC---nền cho \texttt{sid:1000011} và \texttt{sid:1000012}.

\paragraph{Luật \texttt{sid:1000011}.}
\begin{footnotesize}
\begin{verbatim}
alert tcp any any -> $HOME_NET 102 (msg:"ICS S7COMM Write Var to DB1.DBW2 setpoint";
  flow:to_server,established; content:"|03 00|"; offset:0; depth:2;
  content:"|02 F0 80 32 01|"; offset:4; depth:5; content:"|05|";
  offset:17; depth:1; content:"|12 0A 10|"; offset:19; depth:3;
  content:"|00 01 84 00 00 10|"; offset:25; depth:6; sid:1000011; rev:1;
  metadata:service iso-tsap, protocol s7comm, attack stuxnet_mitm_sim; priority:1;)
\end{verbatim}
\end{footnotesize}

\noindent\textbf{Phân tích.}
\texttt{|12 0A 10|}: Any-type; \texttt{|00 01 84 00 00 10|}: DB1, area DB, bit offset \texttt{0x10} (byte~2, \texttt{DBW2}). Khớp địa chỉ \texttt{iSetpoint} trong mô phỏng.

\paragraph{Luật \texttt{sid:1000012}.}
\begin{footnotesize}
\begin{verbatim}
alert tcp any any -> $HOME_NET 102 (msg:"ICS S7COMM Write Var DB1.DBW2 setpoint to 2000";
  flow:to_server,established; content:"|03 00|"; offset:0; depth:2;
  content:"|02 F0 80 32 01|"; offset:4; depth:5; content:"|05|";
  offset:17; depth:1; content:"|12 0A 10|"; offset:19; depth:3;
  content:"|00 01 84 00 00 10|"; offset:25; depth:6;
  content:"|00 04 00 10 07 D0|"; offset:31; depth:6; sid:1000012; rev:1;
  metadata:service iso-tsap, protocol s7comm, attack stuxnet_mitm_sim; priority:1;)
\end{verbatim}
\end{footnotesize}

\noindent\textbf{Phân tích.}
\texttt{|00 04 00 10 07 D0|}: transport size WORD, giá trị \texttt{0x07d0} (2000). Khớp payload tấn công ghi setpoint trong lab. Chuỗi \texttt{sid:1000010}$\rightarrow$\texttt{sid:1000011}$\rightarrow$\texttt{sid:1000012} tăng dần độ đặc hiệu.

\paragraph{Luật \texttt{sid:1000013}.}
\begin{footnotesize}
\begin{verbatim}
alert tcp any 102 -> $HOME_NET any (msg:"ICS S7COMM suspicious Ack_Data response
  to HMI/engineering station"; flow:from_server,established;
  content:"|03 00|"; offset:0; depth:2; content:"|02 F0 80 32 03|";
  offset:4; depth:5; sid:1000013; rev:1;
  metadata:service iso-tsap, protocol s7comm, attack possible_mitm_observation;
  threshold:type both, track by_src, count 60, seconds 10; priority:3;)
\end{verbatim}
\end{footnotesize}

\noindent\textbf{Phân tích.}
Chiều PLC$\rightarrow$HMI, \texttt{ROSCTR=Ack\_Data}. Chữ ký không chứng minh sửa payload MITM; \texttt{threshold} (60 gói/10\,s) ghi nhận mật độ phản hồi cao. Cần tương quan với \texttt{sid:1000012} và phát hiện ARP/L2 (ngoài tập luật S7comm).

\subsection{Kịch bản chưa có luật}
\label{sec:rules-pending}

Mục~\ref{sec:dos_plc} và Mục~\ref{sec:malformed_plc} chưa có PCAP và \texttt{sid} trong \texttt{s7comm.rules}. Khi xác định được byte đặc trưng, bổ sung nhóm \texttt{100004x} theo cùng schema offset ở Mục~\ref{sec:rule-foundation}.

\section{Tổng hợp}
\label{sec:rules-summary}

\begin{table}[H]
    \centering
    \caption{Ánh xạ kịch bản Chương~3 và tập \texttt{sid}}
    \label{tab:rules-master}
    \small
    \begin{tabular}{|p{3.4cm}|l|p{6.5cm}|}
        \hline
        \textbf{Kịch bản} & \textbf{\texttt{sid}} & \textbf{Chức năng phát hiện} \\
        \hline
        Trinh sát & 1000001--1000005 & Setup, SZL, List blocks \\
        \hline
        Start/Stop replay & 1000030--1000033 & \texttt{0x28}/\texttt{0x29}, \texttt{P\_PROGRAM} \\
        \hline
        Upload/Download & 1000020--1000026 & \texttt{0x1a}--\texttt{0x1f}, hai chiều \texttt{0x1b} \\
        \hline
        Stuxnet MITM & 1000010--1000013 & Write \texttt{DB1.DBW2}, quan sát Ack\_Data \\
        \hline
        DoS / Malformed & --- & Chưa triển khai \\
        \hline
    \end{tabular}
\end{table}

\textbf{Hạn chế phương pháp chữ ký:} (i)~chỉ khớp biến thể đã mô tả; (ii)~offset trên payload TCP giả định PDU không bị phân mảnh; (iii)~lưu lượng HMI hợp lệ có thể kích hoạt \texttt{sid:1000001}, \texttt{sid:1000010}; (iv)~\texttt{sid:1000013} không chứng minh MITM độc lập. Các hạn chế này là tiền đề đánh giá thực nghiệm ở chương sau.

\end{document}
